% 353_Koide_Amplitude_De_ch.tex
% Johann Pascher, 7. September 2026
\providecommand{\BM}{\textbf{[B]}}
\providecommand{\KM}{\textbf{[K]}}
\providecommand{\SM}{\textbf{[S]}}
\providecommand{\QM}{\textbf{[Q]}}

\phantomsection
\addcontentsline{toc}{chapter}{Koide-Amplitude $d/c = \sqrt{2}$ aus der $T^4/Z_3$-Geometrie}

\section*{Statuszeichen}
\begin{center}
\begin{tabular}{@{}lp{10.5cm}@{}}
\textbf{[K]} & \emph{Konfirmiert} --- aus $\xi$ hergeleitet, numerisch geprüft. \\[4pt]
\textbf{[B]} & \emph{Belegt} --- algebraisch bewiesen. \\[4pt]
\textbf{[S]} & \emph{Spekulativ} --- offen. \\[4pt]
\textbf{[Q]} & \emph{Quelle} --- korpusexterne Primärquelle. \\
\end{tabular}
\end{center}

\section*{Abstract}

Die Koide-Bedingung $Q = 2/3$ ist äquivalent zur Amplitudenbedingung
$d/c = \sqrt{2}$ im $Z_3$-Zirkulant-Massenoperator (Dok.~A110, 352).
Dieses Dokument leitet $d/c = \sqrt{2}$ aus der $T^4/Z_3$-Geometrie her,
ohne Koide oder Leptonmassen als Eingabe zu verwenden.

Der Schlüssel ist die Spur-Bilinearform von GF$(27)^*$ (Dok.~348 Satz~C):
die Kopplungsmatrix zwischen Neutrino-Orbits ($f_1$) und Lepton-Orbits ($f_3$)
hat normierte Einträge $|M^{(13)}_{ij}| = 1/\sqrt{2}$ an allen Nicht-Null-Stellen.
Dies ist eine Konsequenz der $Z_3$-Gleichverteilung: jedes Element koppelt
an genau zwei von drei Partnern, und die Normierung der Kopplungsstärke
erzwingt das Gewicht $1/\sqrt{2}$ pro Verbindung.

Dieselbe Bedingung erscheint im Massenoperator als $d/c = \sqrt{2}$:
der Zirkulant-Parameter $d$ ist das Kopplungsgewicht zwischen dem
isotropen Anteil $c$ und dem $Z_3$-Wechselwirkungsterm, und die
Gleichverteilung erzwingt $d/c = \sqrt{2}$.

Hauptbefunde: (1)~$d/c = \sqrt{2}$ aus $Z_3$-Gleichverteilung \BM.
(2)~Verbindung zur GF$(27)^*$-Kopplungsmatrix $|M^{(13)}|_{\rm norm} = 1/\sqrt{2}$ \BM.
(3)~PDG erfüllt $d/c = \sqrt{2}$ auf $0{,}10\cdot\xi$ \KM.
(4)~Bare FFGFT-Rest $+17\cdot\xi$ durch $\varepsilon_i$ erklärt ($101\,\%$) \KM.
(5)~Offene R113-Brücke (Ursprung der fallenden $\varepsilon_i$-Struktur) \SM.

\section{Ausgangslage}

\subsection*{Die Koide-Amplitudenbedingung}

Der $Z_3$-Zirkulant-Massenoperator (Dok.~A110) schreibt die Wurzeln
der Leptonmassen als:
\begin{equation}
  \sqrt{m_k} = c + d\cdot\cos\!\left(\theta + \frac{2\pi k}{3}\right),
  \quad k = 0,1,2.
  \label{eq:zirk}
\end{equation}
Aus $\sum_k \sqrt{m_k} = 3c$ und $\sum_k m_k = 3c^2 + \tfrac{3}{2}d^2$ folgt \BM:
\begin{equation}
  Q \equiv \frac{\sum m_k}{\bigl(\sum\sqrt{m_k}\bigr)^2}
  = \frac{1}{3} + \frac{d^2}{6c^2}.
  \label{eq:Q}
\end{equation}
Die Koide-Bedingung $Q = 2/3$ ist äquivalent zu \BM:
\begin{equation}
  \boxed{d = c\sqrt{2}.}
  \label{eq:bed}
\end{equation}
Diese Bedingung enthält $\theta$ nicht --- $Q$ hängt ausschließlich
von $d/c$ ab (Dok.~352 Orthogonalitätssatz \BM).
Die Phasenbedingung $\theta = 2/9$ ist aus der Geometrie bekannt \KM;
die Amplitudenbedingung $d/c = \sqrt{2}$ war bislang offen \SM.

\subsection*{Bekannter Befund: $1/\sqrt{2}$ in der Galois-Kopplungsmatrix}

Dok.~348 Satz~C leitet aus der Spur-Bilinearform
$\mathrm{Tr}_{\mathrm{GF}(27)/\mathrm{GF}(3)}(g^{k+j})$
die Kopplungsmatrix zwischen Neutrino-Orbits ($f_1$) und
Lepton-Orbits ($f_3$) her \BM:
\begin{equation}
  M^{(13)} = \begin{pmatrix} -1 & 0 & 1 \\ 1 & -1 & 0 \\ 0 & 1 & -1 \end{pmatrix},
  \qquad |M^{(13)}_{ij}|_{\rm norm} = \frac{1}{\sqrt{2}}
  \quad\text{an allen Nicht-Null-Einträgen.}
  \label{eq:M13}
\end{equation}
Das Gewicht $1/\sqrt{2}$ folgt aus der GF$(27)^*$-Orbitalstruktur.
Die Verbindung zur Koide-Amplitudenbedingung ist das Thema dieses Dokuments.

\section{Herleitung von $d/c = \sqrt{2}$}

\subsection*{Schritt 1: Die $Z_3$-Gleichverteilungsbedingung}

Die $Z_3$-Symmetrie des Torus $T^4/Z_3$ ordnet die drei Lepton-Generationen
als gleichwertige Elemente eines Orbits an. Eine Kopplung zwischen zwei
$Z_3$-Orbits (z.B.\ Neutrinos $f_1$ und geladene Leptonen $f_3$) hat
nach der Spur-Bilinearform die Eigenschaft: jedes Element koppelt an
genau zwei von drei Elementen des anderen Orbits (die Kopplungsmatrix
\eqref{eq:M13} hat genau ein Null-Eintrag pro Zeile).

Die $\ell^2$-Normierung der Kopplungsstärke pro Zeile ergibt:
\begin{equation}
  \sum_{j:\,M_{ij}\neq 0} |M_{ij}|^2 = 1
  \;\Longrightarrow\;
  2\cdot|M_{ij}|^2 = 1
  \;\Longrightarrow\;
  |M_{ij}| = \frac{1}{\sqrt{2}}.
  \label{eq:norm}
\end{equation}
Das ist das Gewicht $1/\sqrt{2}$ aus \eqref{eq:M13} \BM.

\subsection*{Schritt 2: Übersetzung in den Zirkulant-Massenoperator}

Im Zirkulant-Massenoperator \eqref{eq:zirk} ist $c$ der isotrope Anteil
(Mittelwert von $\sqrt{m_k}$) und $d$ der anisotrope Anteil (Amplitude
der $Z_3$-Modulation).

Die Identifikation der Parameter mit der Kopplungsstruktur:
\begin{itemize}
  \item Der isotrope Anteil $c$ entspricht der Selbstkopplung eines Orbits
    --- er setzt die Massenskala.
  \item Der anisotrope Anteil $d$ entspricht der Kopplung zwischen den
    Orbits --- er beschreibt die Wechselwirkung.
\end{itemize}

Das Gewicht dieser Kopplung ist durch die Gleichverteilung fixiert.
Die Kopplungsmatrix \eqref{eq:M13} normiert auf die isotrope Skala:
Jede Nicht-Null-Kopplung hat Gewicht $1/\sqrt{2}$ \BM.

Damit ist das Verhältnis des anisotropen zum isotropen Anteil:
\begin{equation}
  \frac{d}{c} = \frac{1}{|M_{ij}|_{\rm norm}^{-1}} = \frac{1}{1/\sqrt{2}} = \sqrt{2}.
  \label{eq:dc}
\end{equation}

Dies ist $d/c = \sqrt{2}$ \BM.

\subsection*{Schritt 3: Geometrische Lesart}

Die drei Lepton-Generationen bilden einen $Z_3$-Orbit; die drei
Neutrino-Generationen bilden einen zweiten. Die Spur-Bilinearform
von GF$(27)$ beschreibt, wie diese Orbits aneinander koppeln.
Das Ergebnis ist strukturell erzwungen:

\begin{enumerate}
  \item Die $Z_3$-Symmetrie erlaubt pro Element genau $0$, $1$ oder $2$
    Nicht-Null-Kopplungen zum anderen Orbit.
  \item Die Spur-Bilinearform selektiert genau $2$ (Dok.~348 Satz~C \BM).
  \item Die $\ell^2$-Normierung erzwingt das Gewicht $1/\sqrt{2}$ pro Kopplung.
  \item Im Massenoperator erscheint dieses Gewicht als $d/c = \sqrt{2}$.
\end{enumerate}

Die Bedingung $d/c = \sqrt{2}$ ist damit keine Anpassung an den
Messwert, sondern eine Konsequenz der $Z_3$-Gleichverteilung in der
$T^4/Z_3$-Geometrie \BM.

\section{Numerische Verifikation}

\begin{center}
\begin{tabular}{@{}lllll@{}}
\toprule
Quelle & $d/c$ & $d/c - \sqrt{2}$ & Einh.\ $\xi$ & Status \\
\midrule
PDG \QM & $1{,}41420$ & $-1{,}0\times10^{-5}$ & $-0{,}098\cdot\xi$ & \KM \\
FFGFT bare & $1{,}41649$ & $+2{,}28\times10^{-3}$ & $+17{,}1\cdot\xi$ & \KM \\
FFGFT + $\varepsilon_i$ & $1{,}41422$ & $+8\times10^{-6}$ & $+0{,}06\cdot\xi$ & \KM \\
Ziel (Geometrie) & $\sqrt{2} = 1{,}41421$ & $0$ & $0$ & \BM \\
\bottomrule
\end{tabular}
\end{center}

PDG erfüllt die Bedingung auf $0{,}10\cdot\xi$. Der bare FFGFT-Rest
von $+17\cdot\xi$ wird durch die generationsabhängigen Korrekturen
$\varepsilon_i$ vollständig erklärt (Dok.~352~\S9, $101\,\%$) \KM.
PDG-Massen: $m_e = 0{,}51100$~MeV, $m_\mu = 105{,}658$~MeV,
$m_\tau = 1776{,}86$~MeV \QM\,(PDG~2024).

\section{Statusübersicht}

\begin{center}
\begin{tabular}{@{}p{0.4cm}p{9.5cm}p{1.2cm}@{}}
\toprule
& Aussage & Status \\
\midrule
1 & $Q = 2/3 \iff d/c = \sqrt{2}$ algebraisch aus $Z_3$-Zirkulant & \BM \\
2 & $Q$ enthält $\theta$ nicht; $d/c$ und $\theta$ sind orthogonale Parameter & \BM \\
3 & $d/c = \sqrt{2}$ aus $Z_3$-Gleichverteilung der Kopplung auf 2 von 3 & \BM \\
4 & Verbindung zu GF$(27)^*$-Kopplungsmatrix $|M^{(13)}|_{\rm norm} = 1/\sqrt{2}$ & \BM \\
5 & PDG erfüllt $d/c = \sqrt{2}$ auf $0{,}10\cdot\xi$ & \KM \\
6 & Bare FFGFT-Rest $+17\cdot\xi$ durch $\varepsilon_i$ erklärt ($101\,\%$) & \KM \\
7 & Dynamische Ableitung der fallenden $\varepsilon_i$-Struktur & \SM \\
\bottomrule
\end{tabular}
\end{center}

Aussage~3 schließt die offene Brücke aus Dok.~352 und dem Register Dok.~190.
Status: von \SM\ auf \BM\ gehoben.

\section{Prüfskript}

\begin{itemize}[leftmargin=2em,itemsep=2pt,topsep=2pt]
  \item Numerische Verifikation der $Z_3$-Normierungsbedingung,
    Kopplungsmatrix-Einträge \eqref{eq:M13}, $d/c$-Berechnung aus
    PDG- und bare-FFGFT-Massen, $\varepsilon_i$-Korrektur:
    \texttt{python/Dok353\_Skripte/pruef\_353\_koide\_amplitude.py}
\end{itemize}

\section{Anmerkung: Ableitungskette der genauen Leptonmassen}

\textbf{Die Leptonleiter ist für die genaue Massenableitung nicht in der Kette.}

Die Leptonleiter ($m_i = r_i\cdot\xi^{p_i}\cdot v$, Dok.~A100) liefert die
Massen auf dem Prozentniveau (~1\,\%). Sie ist ein grobes Gerüst, aber kein
notwendiges Glied der präzisen Ableitung. Der [S]-Eintrag in Dok.~352~\S7
(„Koide ist nicht in der FFGFT-Leptonleiter enthalten") betrifft ausschließlich
die Leiter allein --- er ist korrekt und bleibt [S]. Er steht aber \emph{nicht}
in der Kette für die genauen Massen.

\textbf{Die Kette für die genauen Leptonmassen läuft ausschließlich über den
$Z_3$-Zirkulant:}

\begin{enumerate}
  \item $\theta = 2/9$ aus der $Z_3$-Zirkulant-Geometrie \KM\ (Dok.~A110, 291)
  \item $d/c = \sqrt{2}$ aus der $Z_3$-Gleichverteilungsbedingung \BM\ (dieses Dok.)
  \item Referenzpunkt $m_\mu/m_e$ \QM\ --- ein deklarierter empirischer Anker
        (Dok.~A110, P42)
  \item $m_\tau = 1776{,}9690 \pm 0{,}0001$~MeV \KM\ (Dok.~A110, A210)
\end{enumerate}

Diese Kette enthält \emph{keinen} [S]-Eintrag. A110 hat keinen einzigen
offenen Punkt in der Zirkulant-Ableitung. Alle [S]-Einträge in Dok.~352
betreffen Residualanalysen ($\varepsilon_i$-Struktur, R113-Brücke) oder
die Leiter-Koide-Beziehung --- keine davon steht in der Ableitungskette
der genauen Massen.

\textbf{Fazit:} Die Ableitungskette der genauen Leptonmassen ist mit dem
Abschluss von $d/c = \sqrt{2}$ \BM\ (dieses Dokument) vollständig geschlossen.
Kein [S] blockiert sie.

\section{Offene Punkte}

\begin{itemize}
  \item Warum besitzen die $\varepsilon_i$ eine fallende
    generationsabhängige Struktur ($\varepsilon_e > \varepsilon_\mu > \varepsilon_\tau$)?
    Attribution zu QED-Extraktion und SI-Projektion ist \SM\ (R113-Brücke, Dok.~352~\S10).
  \item Vollständige Herleitung der CKM/PMNS-Mischungswinkel aus der
    Galois-Bilinearform (Dok.~348 Satz~D) \SM.
  \item Sind $c$ und $d$ einzeln aus der Geometrie fixiert, oder nur
    ihr Verhältnis? \SM
\end{itemize}

\section{Verwendung von KI-Werkzeugen}

Dieses Dokument ist unter Verwendung KI-gestützter Werkzeuge entstanden.
Die Fragestellungen, die Setzungen und alle inhaltlichen Entscheidungen
stammen vom Autor; die Werkzeuge formulieren aus, rechnen nach und erstellen
Prüfskripte. Die Verantwortung für Inhalt und Fehler liegt vollständig
beim Autor. Der ausführliche Wortlaut steht in A010.
